% =============================================================================
% Бүлэг 2: Арга зүй ба дэвшүүлсэн санаа
% =============================================================================

\chapter{Арга зүй ба дэвшүүлсэн санаа}
\label{ch:method}
\pagecolor{white}

Энэ бүлэгт санал болгож буй системийн архитектур, түүний бүрэлдэхүүн болох SQL үүсгэх загвар ба орчуулгын загварын онолын
үндэс, туршилтын арга зүйг танилцуулна.

\section{Санал болгож буй загвар ба архитектур}
\label{sec:arch}

Системийн гол санаа нь Text2SQL-ийг ганц загварын даалгавар биш, \textbf{асааж унтраах боломжтой алхмуудын дамжлага}
гэж үзэх явдал юм. Алхам бүр туршилтын тохиргооны (\tsql{config/experiments.yaml}) нэг талбараар удирдагддаг тул нөлөөг
нь бусдаас тусад нь хэмжиж болно. Мөн «тогтмол мэдлэгийг загварын жинд, өөрчлөгддөг мэдлэгийг RAG-д» гэсэн зарчим
баримтлав: SQL бичих чадвар, хэлний ойлголт зэрэг тогтвортой зүйлийг сургалтаар, байгууллагын схем, нэр томьёо, утга
зэрэг өөрчлөгдөж болох зүйлийг хайлтаар өгнө. Дамжлагыг Зураг~\ref{fig:pipeline}-т үзүүлэв.

\begin{figure}[htbp]
\centering
\begin{tikzpicture}[node distance=0.45cm and 0.42cm, every node/.style={font=\scriptsize},
  box/.style={draw, rounded corners, align=center, minimum height=1.0cm, text width=2.15cm, fill=gray!8},
  hl/.style={box, fill=orange!12},
  db/.style={draw, cylinder, shape border rotate=90, aspect=0.15, align=center, minimum height=1.0cm,
             minimum width=3.0cm, fill=blue!5}]
  \node[box] (q) {Асуулт\\(кирилл / латин)};
  \node[box, right=of q] (ic) {Шалгах\\(intent, дүрэм)};
  \node[hl, right=of ic] (tr) {Орчуулах\\(Qwen3.5-9B,\\маск)};
  \node[box, right=of tr] (ret) {Мэдлэг хайх\\(BM25 + BGE-M3)};
  \node[hl, below=1.1cm of ret] (gen) {SQL үүсгэх\\Arctic-R1-7B\\+ LoRA (D2)};
  \node[box, left=of gen] (fix) {Хөрвүүлэх, засах\\(sqlglot, ≤2)};
  \node[box, left=of fix] (vote) {8 хувилбар,\\санал хураах};
  \node[box, left=of vote] (out) {Хариулт /\\татгалзах /\\тодруулах};
  \node[db, above=0.55cm of ret] (kn) {Мэдлэгийн сан};
  \node[db, below=0.6cm of fix] (pg) {PostgreSQL (t2s)};
  \draw[->] (q) -- (ic); \draw[->] (ic) -- (tr); \draw[->] (tr) -- (ret); \draw[->] (ret) -- (gen);
  \draw[->] (gen) -- (fix); \draw[->] (fix) -- (vote); \draw[->] (vote) -- (out);
  \draw[<->, dashed] (ret) -- (kn); \draw[<->, dashed] (fix) -- (pg);
  \draw[->] (ic.south) |- ++(-1.2,-0.55) -| node[pos=0.25, above, font=\tiny]{татгалзах} (out.north);
\end{tikzpicture}
\caption{Санал болгож буй дамжлага (улбар шар: загвар ашигладаг алхам)}
\label{fig:pipeline}
\end{figure}

\section{Системийн архитектурын дэлгэрэнгүй тайлбар}
\label{sec:components}

Дамжлага долоон алхамтай (хэрэгжүүлэлтийг Бүлэг~\ref{ch:system}-д дэлгэрэнгүй тайлбарлана):
\begin{enumerate}
  \item \textbf{Шалгах.} Латин галигийг кирилл рүү хөрвүүлж, хувийн мэдээлэл ба өгөгдөл өөрчлөх хүсэлтээс дүрмээр татгалзана.
  \item \textbf{Орчуулах.} Монгол асуултыг англи руу орчуулна (translate-test); DB-д хадгалагдсан утгыг маскаар хамгаална.
  \item \textbf{Мэдлэг хайх.} Схемийн тайлбар, утга, нэр томьёо, жишээг монгол асуултаар хайж prompt-д оруулна.
  \item \textbf{SQL үүсгэх.} Arctic-Text2SQL-R1-7B (сонголтоор LoRA adapter) SQLite-аар бичиж, sqlglot-оор PostgreSQL руу хөрвүүлнэ.
  \item \textbf{Хамгаалах, засах.} SQL-ийг шалгаж, жижиг алдааг дүрмээр, гүйцэтгэлийн алдааг мессежээр нь засуулна.
  \item \textbf{Санал хураах.} 8 хувилбарыг гүйцэтгэж үр дүнгээр нь санал хураана; олонхын хувь нь итгэлийн хэмжүүр болно.
  \item \textbf{Хариулах.} Итгэл хангалттай бол хариулна, бага бол тодруулга асууна, хариулах боломжгүй бол татгалзана.
\end{enumerate}

\subsection{Системийн шаардлага}
\label{sec:requirements}

Сэдэв батлах даалгаврын зорилт ба ББСБ-ийн хэрэглээнээс системийн шаардлагыг тодорхойлов (Хүснэгт~\ref{tab:req}).

\begin{table}[H]
\centering\small
\caption{Функциональ ба функциональ бус шаардлага}
\label{tab:req}
\begin{tabularx}{\textwidth}{p{1.0cm}X p{3.6cm}}
\toprule
№ & Шаардлага & Хэрэгжүүлсэн алхам \\
\midrule
\multicolumn{3}{l}{\textit{Функциональ шаардлага}} \\
Ф1 & Кирилл болон латин галигаар бичсэн монгол асуулт хүлээн авах & Шалгах, хөрвүүлэх \\
Ф2 & Асуултад хамааралтай хүснэгт, утга, нэр томьёо, жишээг олж загварт өгөх & Мэдлэг хайх \\
Ф3 & PostgreSQL дээр ажиллах SQL үүсгэж гүйцэтгэн хариултыг хүснэгтээр буцаах & SQL үүсгэх, гүйцэтгэх \\
Ф4 & Хувийн мэдээлэл, өгөгдөл өөрчлөх хүсэлтээс татгалзах & Шалгах, хамгаалах \\
Ф5 & Хоёрдмол асуултад тодруулга асуух & Санал хураах (итгэл) \\
Ф6 & Хариулт бүрт ашигласан SQL ба хүснэгтийг харуулах & Хариулах, Studio \\
\multicolumn{3}{l}{\textit{Функциональ бус шаардлага}} \\
ФБ1 & Өгөгдөл байгууллагаас гадагш гарахгүй (нээлттэй эхийн загвар) & Нэг GPU, локал vLLM \\
ФБ2 & Зөвхөн унших эрх, салбарын хүрээ (RLS), гүйцэтгэлийн хугацааны хязгаар & DB-ийн эрх, timeout \\
ФБ3 & Туршилт давтагдах боломжтой (seed, hash, тохиргоо) & Run-ийн мета файл \\
ФБ4 & Нэг асуултад 30 секундээс бага хугацаа (нэг үүсгэлттэй горим) & I3t: ~20--30 с (T4) \\
\bottomrule
\end{tabularx}
\end{table}

Нэг асуулт системээр дамжих дарааллыг Зураг~\ref{fig:sequence}-т үзүүлэв.

\begin{figure}[htbp]
\centering
\begin{tikzpicture}[every node/.style={font=\scriptsize}, x=0.98cm, y=-0.8cm, lbl/.style={fill=white, inner sep=1pt, yshift=6pt}]
  \foreach \x/\n in {0/Хэрэглэгч, 2.4/Систем, 4.8/Орчуулагч, 7.2/Мэдлэгийн сан, 9.6/SQL загвар, 12.0/PostgreSQL} {
    \node[draw, fill=gray!10, minimum width=2.0cm, align=center] at (\x,0) {\n};
    \draw[dashed, gray] (\x,0.6) -- (\x,11.4);
  }
  \draw[->] (0,1.2) -- node[lbl]{монгол асуулт} (2.4,1.2);
  \draw[->] (2.4,2.1) -- node[lbl]{асуулт (маскласан)} (4.8,2.1);
  \draw[<-] (2.4,3.0) -- node[lbl]{англи орчуулга} (4.8,3.0);
  \draw[->] (2.4,3.9) -- node[lbl, pos=0.7]{монгол асуултаар хайх} (7.2,3.9);
  \draw[<-] (2.4,4.8) -- node[lbl, pos=0.7]{хүснэгт, утга, нэр томьёо, жишээ} (7.2,4.8);
  \draw[->] (2.4,5.7) -- node[lbl, pos=0.75]{prompt (англи асуулт + мэдлэг)} (9.6,5.7);
  \draw[<-] (2.4,6.6) -- node[lbl, pos=0.75]{8 SQL хувилбар} (9.6,6.6);
  \draw[->] (2.4,7.5) -- node[lbl, pos=0.8]{хувилбар бүрийг гүйцэтгэх (read-only)} (12.0,7.5);
  \draw[<-] (2.4,8.4) -- node[lbl, pos=0.8]{үр дүн / алдаа} (12.0,8.4);
  \draw[->] (2.4,9.3) -- ++(0.6,0) -- ++(0,0.6) node[right, align=left]{санал хураах, итгэл} -- ++(-0.6,0);
  \draw[<-] (0,10.8) -- node[lbl]{хариулт / тодруулга} (2.4,10.8);
\end{tikzpicture}
\caption{Нэг асуултын дарааллын диаграм (E2 тохиргоо)}
\label{fig:sequence}
\end{figure}

\subsection{Туршилтын арга зүй}
\label{sec:exp-method}

\textbf{Хуваалт ба шийдвэрийн дүрэм.} Бүх шийдвэрийг dev (300) дээр гаргаж, test-ийг (200) шилдэг тохиргоонд л
ажиллуулна. Dev ба test нь загварын гэр бүлээр давхцалгүй хуваагдсан (§\ref{sec:evalset}).

\textbf{Шударга байдлын дүрэм.} (1) Dev-ийн асуулт тус бүрд тааруулахгүй, зөвхөн бүх тохиргоонд ижил үйлчлэх ерөнхий
засвар хийнэ; (2) test-ийн асуулт жишээний санд хэзээ ч орохгүй, dev-ийн асуулт өөрийн гэр бүлийн жишээг татахгүй;
(3) шинэ жишээ бичих agent-уудад үнэлгээний санг уншихыг хориглож, үнэлгээний асуулттай 3-gram cosine ≥ 0.8 бол хасна;
(4) dev-ийг харж гаргасан дүрэм, босгыг (E4-ийн түлхүүр үг, E3c-ийн 0.5) ил тод тэмдэглэнэ.

\textbf{Хэмжүүр.} TrustSQL маягаар \cite{lee_2024_trustsql}: хариулах боломжтой асуулт --- гүйцэтгэлийн үр дүн gold-той
таарвал зөв; хариулах боломжгүй --- татгалзсан бол зөв; хоёрдмол --- тодруулга асуусан эсвэл аль нэг тайлбартай таарвал
зөв. Гүйцэтгэлийн нарийвчлал:
\begin{equation}
\mathrm{EX} = \frac{1}{N}\sum_{j=1}^{N} \mathbf{1}\!\left[R(\hat{s}_j, D) = R(s^{\ast}_j, D)\right],
\label{eq:ex}
\end{equation}
энд $R(s, D)$ нь $s$-ийг $D$ сан дээр гүйцэтгэсэн үр дүн (мөрийн multiset, тоог 0.01 хүлцэлтэй харьцуулна). \textbf{Нийт
зөв} (/300, /200) ба \textbf{хариулах боломжтой EX} (/224, /151)-ийг үргэлж хамт тайлагнана: тодруулга нийт хувийг өсгөж,
хариулах боломжтой EX-ийг бууруулж болдог. Итгэх завсрыг bootstrap-аар, хос тохиргооны ялгааг exact McNemar тестээр
\cite{mcnemar_1947} тооцно.

\section{SQL үүсгэх загвар}
\label{sec:sql-model}

\textbf{Arctic-Text2SQL-R1-7B} нь Qwen2.5-Coder-7B-Instruct-д \cite{hui_2024_qwen25_coder} суурилсан, OmniSQL-оос
эхлүүлж GRPO хийсэн загвар \cite{yao_2025_arctic_text2sql_r1, snowflake_2025_arctic_card}. GRPO нь нэг асуултын $G$
хариуны reward-ийг бүлгийн дунджаар хэвийн болгож бодлогыг шинэчилдэг \cite{shao_2024_deepseekmath}. RL-ээр сургасан
загвар өөр prompt хэлбэрт муудаж болох тул түүний сургалтын протоколыг яг мөрдөв: system нь «data science expert», user нь
Database Engine / Database Schema / Question / Instructions / Output Format хэсэгтэй, хариуг «Let me solve this step by
step. <think>» гэж эхлүүлж, SQL-ийг \tsql{<answer>} доторх код блокоос авна.

\begin{table}[H]
\centering\small
\caption{SQL үүсгэх загварын үзүүлэлт}
\label{tab:arctic}
\begin{tabularx}{\textwidth}{p{4.2cm}X}
\toprule
Үзүүлэлт & Утга \\
\midrule
Суурь загвар & Qwen2.5-Coder-7B-Instruct (decoder-only transformer, ~7B параметр) \\
Сургалт & OmniSQL-ийн SFT → гүйцэтгэлийн reward-тай GRPO \\
Жишиг үр дүн & BIRD test 68.5\% EX \cite{yao_2025_arctic_text2sql_r1} \\
Prompt-ийн диалект & SQLite (сургалтынх) → sqlglot-оор PostgreSQL \cite{mao_2026_sqlglot} \\
Декодчилол & Greedy ($T = 0$); санал хураалтад sampling-аар 8 хувилбар \\
Ажиллуулах орчин & vLLM \cite{kwon_2023_vllm}; T4×2 (fp16) эсвэл A100 (bf16), контекст 32 768 токен \\
\bottomrule
\end{tabularx}
\end{table}

\textbf{LoRA ба QLoRA.} LoRA нь жингийн матрицад бага rank-тай засвар $\Delta W = BA$ ($r \ll \min(d,k)$) нэмж сургадаг тул
7B загварыг нэг GPU дээр сургах боломж олгоно \cite{hu_2022_lora}. QLoRA нь суурь жинг 4 битээр шахаж санах ойг цаашид
багасгана \cite{dettmers_2023_qlora}. Rejection sampling SFT нь загвараар олон хариу үүсгээд гүйцэтгэлийн үр дүн нь
gold-той таарсныг нь сургалтад үлдээдэг.

\section{Орчуулгын загвар ба мэдлэг хайлт}
\label{sec:mt-model}

\textbf{Орчуулга.} Arctic загвар англи өгөгдлөөр сургагдсан тул монгол асуултыг англи руу орчуулж өгнө. Одоогийн
орчуулагч нь Qwen3.5-9B (zero-shot, заавартай) бөгөөд нэмж сургах нэр дэвшигч нь 400 гаруй хэлний өгөгдлөөр сургасан
MADLAD-400-3B-MT \cite{kudugunta_2023_madlad} юм. Орчуулгын чанарыг тэмдэгт ба үгийн n-gram-д суурилсан chrF++ хэмжүүрээр
\cite{popovic_2017_chrfpp} үнэлнэ.

\textbf{Мэдлэг хайлт.} Түлхүүр үгийн BM25 \cite{robertson_2009_bm25} нь яг тохирох нэр, кодыг сайн олдог бол олон хэлт
BGE-M3 \cite{chen_2024_bge_m3} утгын төстэй байдлыг барьдаг. Хоёр оноог $\lambda = 0.5$ жингээр хослуулна. Жишээг
DAIL-SQL маягаар асуултын утгыг далдалж сонгоно \cite{gao_2024_dail_sql}.

\begin{table}[H]
\centering\small
\caption{Туслах загваруудын үзүүлэлт}
\label{tab:aux-models}
\begin{tabularx}{\textwidth}{p{3.2cm}p{3.6cm}X}
\toprule
Загвар & Бүтэц, хэмжээ & Үүрэг \\
\midrule
Qwen3.5-9B & Decoder-only, ~9B, 4 бит (T4) & Монгол → англи орчуулга (zero-shot, нэр томьёоны зөвлөмжтэй) \\
MADLAD-400-3B-MT & Encoder--decoder (T5), ~3B & Нэмж сургах орчуулагч (QLoRA, §\ref{sec:mt-train}) \\
BGE-M3 & Олон хэлт embedding & Мэдлэгийн сангийн утгын хайлт \\
\bottomrule
\end{tabularx}
\end{table}

\section{Бүлгийн дүгнэлт}
\label{sec:ch3-summary}

Шалгах, орчуулах, мэдлэг хайх, SQL үүсгэх, засах, санал хураах, хариулах гэсэн долоон алхамтай, алхам бүрийг асааж
унтраах боломжтой дамжлагыг санал болгов. SQL бичигчээр Arctic-R1-7B-г өөрийн протоколоор нь ажиллуулж, англи төвтэй
загварын хэлний зөрүүг орчуулга ба мэдлэг хайлтаар нөхөхөөр төлөвлөв. Туршилтын арга зүй нь dev дээр шийдэж test дээр нэг
удаа батлах, хоёр хэмжүүр ба McNemar тест дээр тулгуурлана.
